% 本文件由账本已入列事件生成；锚点首次呈现仅连接可核的年序与材料限度。
\section{北宋中枢与边地的加密年序}
以下补入按同年并列的政权与地区组织，而不是以一个朝廷的年号吞没其对手。每个可点击事件名保留账本 ID；叙述只将材料明确记录的动作置入年序，并把实施、规模、损失与动机留在可检验的限度内。

\subsection{960年：宋}
《宋史》为这一年留下了可回查的年序入口。

\noindent\textbf{宋。}\eventanchor{SYD-960-NS-001}{八月戊辰朔，御崇元殿，行入閣儀}；\eventanchor{SYD-960-NS-168}{丙辰，南唐主李景、吳越王錢俶遣使以御服、錦綺、金帛來賀}。

\subsection{961年：宋}
《宋史》为这一年留下了可回查的年序入口。

\noindent\textbf{宋。}\eventanchor{SYD-961-NS-002}{壬辰，南唐進謝賜生辰金器、羅綺}；\eventanchor{SYD-961-NS-169}{丙申，遣樞密承旨王仁贍賜南唐禮物}。

\subsection{962年：宋}
《宋史》为这一年留下了可回查的年序入口。

\noindent\textbf{宋。}\eventanchor{SYD-962-NS-003}{八月癸巳，蔡河務綱官王訓等四人坐以糠土雜軍糧，磔於市}；\eventanchor{SYD-962-NS-170}{丙申，寧州大雨雪，溝洫冰}。

\subsection{963年：宋}
《宋史》为这一年留下了可回查的年序入口。

\noindent\textbf{宋。}\eventanchor{SYD-963-NS-004}{八月壬午，殿前都虞候張瓊以陵侮軍校史珪、石漢卿等，為所誣譖，下吏，瓊自殺}；\eventanchor{SYD-963-NS-171}{丙辰，幸新池，賜役夫錢，遂幸玉津園}。

\subsection{964年：宋}
《宋史》为这一年留下了可回查的年序入口。

\noindent\textbf{宋。}\eventanchor{SYD-964-NS-005}{丁巳，治安陵，隧壞，役兵壓死者二百人，命有司瘞恤}；\eventanchor{SYD-964-NS-172}{丁未，詔縣令、簿、尉非公事毋至村落}。

\subsection{965年：宋}
《宋史》为这一年留下了可回查的年序入口。

\noindent\textbf{宋。}\eventanchor{SYD-965-NS-006}{八月戊戌朔，詔籍郡國驍勇兵送闕下}；\eventanchor{SYD-965-NS-173}{丙申，赦蜀，歸俘獲，除管内逋賦，免夏税及沿徵物色之半}。

\subsection{966年：宋}
《宋史》为这一年留下了可回查的年序入口。

\noindent\textbf{宋。}\eventanchor{SYD-966-NS-007}{安國軍節度使羅彥瓌等敗北漢於靜陽，擒其將鹿英}；\eventanchor{SYD-966-NS-174}{八月丁酉，詔除蜀倍息}。

\subsection{967年：宋}
《宋史》为这一年留下了可回查的年序入口。

\noindent\textbf{宋。}\eventanchor{SYD-967-NS-008}{八月甲申，河溢入衛州城，民溺死者數百}；\eventanchor{SYD-967-NS-175}{北漢鴻唐砦招收指揮使樊暉以砦來降}。

\subsection{968年：宋}
《宋史》为这一年留下了可回查的年序入口。

\noindent\textbf{宋。}\eventanchor{SYD-968-NS-009}{北漢穎州砦主胡遇等來降}；\eventanchor{SYD-968-NS-176}{北漢主劉鈞卒，養子繼恩立}。

\subsection{969年：宋}
《宋史》为这一年留下了可回查的年序入口。

\noindent\textbf{宋。}\eventanchor{SYD-969-NS-010}{八月丁亥，詔川峽諸州察民有父母在而別籍異財者，論死}；\eventanchor{SYD-969-NS-177}{賜宰相、樞密使、翰林學士、節度、觀察使襲衣金帶}。

\subsection{970年：宋}
《宋史》为这一年留下了可回查的年序入口。

\noindent\textbf{宋。}\eventanchor{SYD-970-NS-011}{丙辰，殿中丞張顒坐先知潁州政不平，免官}；\eventanchor{SYD-970-NS-178}{丙辰，女直國遣使齎定安國王烈萬華表，獻方物}。

\subsection{971年：宋}
《宋史》为这一年留下了可回查的年序入口。

\noindent\textbf{宋。}\eventanchor{SYD-971-NS-012}{丙申，詔：廣南有賣人男女為奴婢轉傭利者，並放免}；\eventanchor{SYD-971-NS-179}{丙午，罷諸道州縣攝官}。

\subsection{972年：宋}
《宋史》为这一年留下了可回查的年序入口。

\noindent\textbf{宋。}\eventanchor{SYD-972-NS-013}{八月庚寅，高麗國王王昭遣使獻方物}；\eventanchor{SYD-972-NS-180}{丙午，遣使檢視水災田}。

\subsection{973年：宋}
《宋史》为这一年留下了可回查的年序入口。

\noindent\textbf{宋。}\eventanchor{SYD-973-NS-014}{八月乙酉，罷成都府偽蜀嫁裝稅}；\eventanchor{SYD-973-NS-181}{丙申，曹州饑，漕太倉米二萬石振之}。

\subsection{974年：宋}
《宋史》为这一年留下了可回查的年序入口。

\noindent\textbf{宋。}\eventanchor{SYD-974-NS-015}{八月戊寅，吳越國王遣使來朝貢}；\eventanchor{SYD-974-NS-182}{丙戌，又幸迎春苑，登汴隄觀諸軍習戰，遂幸東水門，發戰櫂東下}。

\subsection{975年：宋}
《宋史》为这一年留下了可回查的年序入口。

\noindent\textbf{宋。}\eventanchor{SYD-975-NS-016}{八月乙卯，幸東水磑觀魚，遂幸北園}；\eventanchor{SYD-975-NS-183}{丙子，知池州樊若水敗江南軍於州界}。

\subsection{976年：宋}
《宋史》为这一年留下了可回查的年序入口。

\noindent\textbf{宋。}\eventanchor{SYD-976-NS-017}{汴京新宮成，御正殿坐，令洞開諸門，謂左右曰：「此如我心，少有邪曲，人皆見之}；\eventanchor{SYD-976-NS-184}{帝笑而謂之曰：「朕推赤心於人腹中，寧肯爾耶？」即取鋹酒自飲，別酌以賜鋹}。

\subsection{977年：宋}
《宋史》为这一年留下了可回查的年序入口。

\noindent\textbf{宋。}\eventanchor{SYD-977-NS-018}{八月癸亥，黎州兩林蠻來貢}；\eventanchor{SYD-977-NS-185}{丙寅，禁居官出使者行商賈事}。

\subsection{978年：宋}
《宋史》为这一年留下了可回查的年序入口。

\noindent\textbf{宋。}\eventanchor{SYD-978-NS-019}{八月癸丑，幸南造船務，遂幸玉津園宴射}；\eventanchor{SYD-978-NS-186}{丙辰，禁民自春及秋毋捕獵}。

\subsection{979年：宋}
《宋史》为这一年留下了可回查的年序入口。

\noindent\textbf{宋。}\eventanchor{SYD-979-NS-020}{北漢節度使蔚進盧遂以汾州降}；\eventanchor{SYD-979-NS-187}{丙辰，以中書侍郎、尚書右僕射、同平章事沈倫為東京留守兼判開封府事，宣徽北院使王仁贍為大內都部署，樞密承旨陳從信副之}。

\subsection{980年：宋}
《宋史》为这一年留下了可回查的年序入口。

\noindent\textbf{宋。}\eventanchor{SYD-980-NS-021}{八月甲申，西南蕃主龍瓊琚使其子羅若從並諸州蠻來貢}；\eventanchor{SYD-980-NS-188}{冬十月戊寅，大發兵屯關南及鎮、定州}。

\subsection{981年：宋}
《宋史》为这一年留下了可回查的年序入口。

\noindent\textbf{宋。}\eventanchor{SYD-981-NS-022}{丙辰，置破虜、平戎二軍}；\eventanchor{SYD-981-NS-189}{丙午，置京朝官差遣院，初令中書舍人郭贄等考校課績}。

\subsection{982年：宋}
《宋史》为这一年留下了可回查的年序入口。

\noindent\textbf{宋。}\eventanchor{SYD-982-NS-023}{北陽縣蝗，飛鳥數萬食之盡}；\eventanchor{SYD-982-NS-190}{丙辰，秦王廷美降封涪陵縣公、房州安置}。

\subsection{983年：宋}
《宋史》为这一年留下了可回查的年序入口。

\noindent\textbf{宋。}\eventanchor{SYD-983-NS-024}{八月壬辰，以大水故，釋死罪以下}；\eventanchor{SYD-983-NS-191}{初置水陸路發運于京師}。

\subsection{984年：宋}
《宋史》为这一年留下了可回查的年序入口。

\noindent\textbf{宋。}\eventanchor{SYD-984-NS-025}{丙申，御乾元門，賜京師大酺三日}；\eventanchor{SYD-984-NS-192}{丁卯，涪陵縣公廷美薨，追封涪陵王}。

\subsection{985年：宋}
《宋史》为这一年留下了可回查的年序入口。

\noindent\textbf{宋。}\eventanchor{SYD-985-NS-026}{八月癸酉朔，遣使按問兩浙、荊湖、福建、江南東西路、淮南諸州刑獄，仍察官吏勤惰以聞}；\eventanchor{SYD-985-NS-193}{丙辰，門下侍郎兼刑部尚書、平章事宋琪罷守本官}。

\subsection{986年：宋}
《宋史》为这一年留下了可回查的年序入口。

\noindent\textbf{宋。}\eventanchor{SYD-986-NS-027}{丙申，進圍靈丘，其守將穆超以城降}；\eventanchor{SYD-986-NS-194}{丙子，召曹彬、崔彥進、米信歸闕，命田重進屯定州，潘美還代州}。

\subsection{987年：宋}
《宋史》为这一年留下了可回查的年序入口。

\noindent\textbf{宋。}\eventanchor{SYD-987-NS-028}{八月庚子，免諸州吏所逋京倉米二十六萬七千石}；\eventanchor{SYD-987-NS-195}{丙戌，詔：「應行營將士戰敗潰散者並釋不問，緣邊城堡備禦有勞可紀者所在以聞}。

\subsection{988年：宋}
《宋史》为这一年留下了可回查的年序入口。

\noindent\textbf{宋。}\eventanchor{SYD-988-NS-029}{丙申，禁諸州獻珍禽奇獸}；\eventanchor{SYD-988-NS-196}{除十惡、官吏犯贓至殺人者不赦外，民年七十以上賜爵一級}。

\subsection{989年：宋}
《宋史》为这一年留下了可回查的年序入口。

\noindent\textbf{宋。}\eventanchor{SYD-989-NS-030}{霸、代、洺、雄、莫、深等州，平虜、岢嵐軍，更給復一年}；\eventanchor{SYD-989-NS-197}{二年春正月癸未朔，不受朝，群臣詣閤拜表稱賀}。

\subsection{990年：宋}
《宋史》为这一年留下了可回查的年序入口。

\noindent\textbf{宋。}\eventanchor{SYD-990-NS-031}{丙寅，命殿前副都指揮使戴興為鎮州都部署}；\eventanchor{SYD-990-NS-198}{曹、單二州有蝗，不為災}。

\subsection{991年：宋}
《宋史》为这一年留下了可回查的年序入口。

\noindent\textbf{宋。}\eventanchor{SYD-991-NS-032}{丙辰，左正言謝泌以敢言擢右司諫，賜金紫，錢三十萬}；\eventanchor{SYD-991-NS-199}{大名、河中，絳、濮、陝、曹、濟、同、淄、單、德、徐、晉、輝、磁、博、汝、兗、虢、汾、鄭、亳、慶、許、齊、濱、棣、沂、貝、衛、青、霸等州旱}。

\subsection{992年：宋}
《宋史》为这一年留下了可回查的年序入口。

\noindent\textbf{宋。}\eventanchor{SYD-992-NS-033}{八月戊辰，以秘閤成，賜近臣宴}；\eventanchor{SYD-992-NS-200}{丙午，詔宰相、侍從舉可任轉運使者}。

\subsection{993年：宋}
《宋史》为这一年留下了可回查的年序入口。

\noindent\textbf{宋。}\eventanchor{SYD-993-NS-034}{丙午，命侍從舉任才堪五千戶以上縣令者二人}；\eventanchor{SYD-993-NS-201}{丙戌，置審官院、考課院}。

\subsection{994年：宋}
《宋史》为这一年留下了可回查的年序入口。

\noindent\textbf{宋。}\eventanchor{SYD-994-NS-035}{八月甲申，詔有司講求大射儀注}；\eventanchor{SYD-994-NS-202}{貝州言驍捷卒劫庫兵為亂，推都虞候趙咸雍為帥，轉運使王嗣宗率屯兵擊敗之，擒咸雍，磔於市}。

\subsection{995年：宋}
《宋史》为这一年留下了可回查的年序入口。

\noindent\textbf{宋。}\eventanchor{SYD-995-NS-036}{八月壬辰，詔立壽王元侃為皇太子，改名恆，兼判開封府}；\eventanchor{SYD-995-NS-203}{丙戌，遣使諭李繼遷，授以鄜州節度使，繼遷不奉詔}。

\subsection{996年：宋}
《宋史》为这一年留下了可回查的年序入口。

\noindent\textbf{宋。}\eventanchor{SYD-996-NS-037}{八月辛丑，密州言蝗不為災}；\eventanchor{SYD-996-NS-204}{丙戌，秦、晉諸州地晝夜十二震}。

\subsection{997年：宋}
《宋史》为这一年留下了可回查的年序入口。

\noindent\textbf{宋。}\eventanchor{SYD-997-NS-038}{二月丙申朔，靈州行營破李繼遷}；\eventanchor{SYD-997-NS-205}{癸巳，追班于萬歲殿，宣詔令皇太子柩前即位}。

\subsection{998年：宋}
《宋史》为这一年留下了可回查的年序入口。

\noindent\textbf{宋。}\eventanchor{SYD-998-NS-039}{丙辰，以旱，免開封、二十五州軍田租}；\eventanchor{SYD-998-NS-206}{帝曰：「朕以天下為憂，豈直一方耶？」甲午，詔求直言，避殿減膳}。

\subsection{999年：宋}
《宋史》为这一年留下了可回查的年序入口。

\noindent\textbf{宋。}\eventanchor{SYD-999-NS-040}{詔天下繫囚非十惡、枉法及已殺人者，死以下減一等}；\eventanchor{SYD-999-NS-207}{八月辛亥，御文德殿，文武百官入閤}。

\subsection{1000年：宋}
《宋史》为这一年留下了可回查的年序入口。

\noindent\textbf{宋。}\eventanchor{SYD-1000-NS-041}{八月辛亥，京東水災，遣使安撫}；\eventanchor{SYD-1000-NS-208}{罷緣邊二十三州軍榷酤}。

\subsection{1001年：宋}
《宋史》为这一年留下了可回查的年序入口。

\noindent\textbf{宋。}\eventanchor{SYD-1001-NS-042}{李沆等乞免官，不許}；\eventanchor{SYD-1001-NS-209}{丙辰，審官院引對京朝官，閱殿最而黜陟之}。

\subsection{1002年：宋}
《宋史》为这一年留下了可回查的年序入口。

\noindent\textbf{宋。}\eventanchor{SYD-1002-NS-043}{丙戌，賜深、霸九州民租有差}；\eventanchor{SYD-1002-NS-210}{丙子，沙州曹宗壽遣使入貢，以宗壽為歸義軍節度使}。

\subsection{1003年：宋}
《宋史》为这一年留下了可回查的年序入口。

\noindent\textbf{宋。}\eventanchor{SYD-1003-NS-044}{丙申，出內府繒帛，市穀實邊}；\eventanchor{SYD-1003-NS-211}{丙子，詔：環慶秋田經寇踐傷者，頃賜粟十五斛}。

\subsection{1004年：宋}
《宋史》为这一年留下了可回查的年序入口。

\noindent\textbf{宋。}\eventanchor{SYD-1004-NS-045}{八月，涇原部署言擊萬子軍主族帳，斬首二百餘級}；\eventanchor{SYD-1004-NS-212}{丙戌，遣使撫諭懷、孟、澤、潞、鄭、滑等州，放強壯歸農}。

\subsection{1005年：宋}
《宋史》为这一年留下了可回查的年序入口。

\noindent\textbf{宋。}\eventanchor{SYD-1005-NS-046}{罷江、淮、荊、浙增榷酤錢}；\eventanchor{SYD-1005-NS-213}{丁丑，詔河北轉運使察官屬不任職者以名聞}。

\subsection{1006年：宋}
《宋史》为这一年留下了可回查的年序入口。

\noindent\textbf{宋。}\eventanchor{SYD-1006-NS-047}{丙申，遣使振應天府水災及瘞溺死者}；\eventanchor{SYD-1006-NS-214}{丙寅，大風，遣中使視稼}。

\subsection{1007年：宋}
《宋史》为这一年留下了可回查的年序入口。

\noindent\textbf{宋。}\eventanchor{SYD-1007-NS-048}{八月壬寅，幸大相國寺，遂幸崇文院觀書，賜修書官器幣}；\eventanchor{SYD-1007-NS-215}{丙辰，涇原路言瓦亭砦地震}。

\subsection{1008年：宋}
《宋史》为这一年留下了可回查的年序入口。

\noindent\textbf{宋。}\eventanchor{SYD-1008-NS-049}{八月己丑，上太祖尊諡曰啟運立極英武聖文神德玄功大孝皇帝，太宗曰至仁應道神功聖德文武大明廣孝皇帝}；\eventanchor{SYD-1008-NS-216}{丙申，以王旦為封禪大禮使，馮拯、陳堯叟分掌禮儀使}。

\subsection{1009年：宋}
《宋史》为这一年留下了可回查的年序入口。

\noindent\textbf{宋。}\eventanchor{SYD-1009-NS-050}{庚午，詔：「讀非聖之書及屬辭浮靡者，皆嚴譴之}；\eventanchor{SYD-1009-NS-217}{乙酉，以陝西民饑，遣使巡撫}。

\subsection{1010年：宋}
《宋史》为这一年留下了可回查的年序入口。

\noindent\textbf{宋。}\eventanchor{SYD-1010-NS-051}{八月丁未朔，詔：明年春有事于汾陰，州府長吏勿以修貢助祭煩民}；\eventanchor{SYD-1010-NS-218}{丙寅，詔沙門島流人特給口糧}。

\subsection{1011年：宋}
《宋史》为这一年留下了可回查的年序入口。

\noindent\textbf{宋。}\eventanchor{SYD-1011-NS-052}{丙申，詔以六月六日天書再降日為天貺節}；\eventanchor{SYD-1011-NS-219}{丙寅，遣使安撫江、淮南水災，許便宜從事}。

\subsection{1012年：宋}
《宋史》为这一年留下了可回查的年序入口。

\noindent\textbf{宋。}\eventanchor{SYD-1012-NS-053}{丙戌，詔天慶等節日，民犯罪情輕者釋之}；\eventanchor{SYD-1012-NS-220}{丙寅，詔官吏安撫濱、棣被水農民}。

\subsection{1013年：宋}
《宋史》为这一年留下了可回查的年序入口。

\noindent\textbf{宋。}\eventanchor{SYD-1013-NS-054}{八月庚申，詔來春親謁亳州太清宮}；\eventanchor{SYD-1013-NS-221}{丙午，詔聖像所經郡邑減系囚死罪，流以下釋之}。

\subsection{1014年：宋}
《宋史》为这一年留下了可回查的年序入口。

\noindent\textbf{宋。}\eventanchor{SYD-1014-NS-055}{八月甲寅，置景靈宮使，以向敏中為之}；\eventanchor{SYD-1014-NS-222}{丙辰，詔王欽若等五人各舉京朝、幕職、州縣官詳練刑典、曉時務、任邊寄者二人}。

\subsection{1015年：宋}
《宋史》为这一年留下了可回查的年序入口。

\noindent\textbf{宋。}\eventanchor{SYD-1015-NS-056}{八年春正月壬午朔，謁玉清昭應宮，奉表告尊上玉皇大天帝聖號，奉安刻玉天書于寶符閣，還禦崇德殿受賀，赦天下，非十惡、枉法贓及己殺人者鹹除之}；\eventanchor{SYD-1015-NS-223}{八月己卯，大理少卿閻允恭、開封判官韓允坐枉獄除名}。

\subsection{1016年：宋}
《宋史》为这一年留下了可回查的年序入口。

\noindent\textbf{宋。}\eventanchor{SYD-1016-NS-057}{八月壬申，知秦州曹瑋言伏羌砦蕃部廝雞波與宗哥族連結為亂，以兵夷其族帳}；\eventanchor{SYD-1016-NS-224}{丙辰，開封府祥符縣蝗附草死者數里}。

\subsection{1017年：宋}
《宋史》为这一年留下了可回查的年序入口。

\noindent\textbf{宋。}\eventanchor{SYD-1017-NS-058}{丙寅，命王旦為兗州太極觀奉上冊寶使}；\eventanchor{SYD-1017-NS-225}{丙子，詔京城禁圍草地聽民耕牧}。

\subsection{1018年：宋}
《宋史》为这一年留下了可回查的年序入口。

\noindent\textbf{宋。}\eventanchor{SYD-1018-NS-059}{八月庚寅，群臣請立皇太子，從之}；\eventanchor{SYD-1018-NS-226}{丙辰，以德雍、德文、惟政並為諸州防禦使，允成、允升、允甯並為諸州團練使}。

\subsection{1019年：宋}
《宋史》为这一年留下了可回查的年序入口。

\noindent\textbf{宋。}\eventanchor{SYD-1019-NS-060}{丙寅，禦試禮部貢舉人}；\eventanchor{SYD-1019-NS-227}{二月乙未，河南府地震}。

\subsection{1020年：宋}
《宋史》为这一年留下了可回查的年序入口。

\noindent\textbf{宋。}\eventanchor{SYD-1020-NS-061}{宰臣丁謂請鏤板宣佈}；\eventanchor{SYD-1020-NS-228}{八月，永興軍都巡檢使朱能殺中使叛}。

\subsection{1021年：宋}
《宋史》为这一年留下了可回查的年序入口。

\noindent\textbf{宋。}\eventanchor{SYD-1021-NS-062}{庚午，以孔子四十七世孫聖祐襲封文宣公}；\eventanchor{SYD-1021-NS-229}{癸巳，詔天下死罪降，流以下釋之}。

\subsection{1022年：宋}
《宋史》为这一年留下了可回查的年序入口。

\noindent\textbf{宋。}\eventanchor{SYD-1022-NS-063}{甲辰，制封丁謂為晉國公，馮拯為魏國公，曹利用為韓國公}；\eventanchor{SYD-1022-NS-230}{乾興元年春正月辛未朔，改元}。

\subsection{1023年：宋}
《宋史》为这一年留下了可回查的年序入口。

\noindent\textbf{宋。}\eventanchor{SYD-1023-NS-064}{八月乙未，募民輸芟塞滑州決河}；\eventanchor{SYD-1023-NS-231}{丙子，詔減西京囚罪一等，徒以下釋之}。

\subsection{1024年：宋}
《宋史》为这一年留下了可回查的年序入口。

\noindent\textbf{宋。}\eventanchor{SYD-1024-NS-065}{罷天慶、天祺、天貺、先天、降聖節宮觀然燈}；\eventanchor{SYD-1024-NS-232}{二年春二月庚午，遣內臣收瘞汴口流屍，仍祭奠之}。

\subsection{1025年：宋}
《宋史》为这一年留下了可回查的年序入口。

\noindent\textbf{宋。}\eventanchor{SYD-1025-NS-066}{八月戊午，以忠州鹽井歲增課、夔州奉節巫山縣舊藉民為營田、萬州戶有稅者歲糴其穀，皆為民害，詔悉除之}；\eventanchor{SYD-1025-NS-233}{丙午，詔邊戶為羌所擾者蠲租，復役二年}。

\subsection{1026年：宋}
《宋史》为这一年留下了可回查的年序入口。

\noindent\textbf{宋。}\eventanchor{SYD-1026-NS-067}{八月丁亥，築泰州捍海堰}；\eventanchor{SYD-1026-NS-234}{除舒州太湖等九茶場民逋錢十三萬緡}。

\subsection{1027年：宋}
《宋史》为这一年留下了可回查的年序入口。

\noindent\textbf{宋。}\eventanchor{SYD-1027-NS-068}{罷瓊州歲貢瑇瑁、龜皮、紫貝}；\eventanchor{SYD-1027-NS-235}{丙辰，發丁夫三萬八千、卒二萬一千、緡錢五十萬塞滑州決河}。

\subsection{1028年：宋}
《宋史》为这一年留下了可回查的年序入口。

\noindent\textbf{宋。}\eventanchor{SYD-1028-NS-069}{八月乙丑，詔免河北水災州軍秋稅}；\eventanchor{SYD-1028-NS-236}{丁-\{丑\}-，貸河北流民復業者種食，復是年租賦}。

\subsection{1029年：宋}
《宋史》为这一年留下了可回查的年序入口。

\noindent\textbf{宋。}\eventanchor{SYD-1029-NS-070}{丙辰，降利用為左千牛衛上將軍}；\eventanchor{SYD-1029-NS-237}{丙寅，張士遜罷，以呂夷簡同中書門下平章事、集賢殿大學士}。

\subsection{1030年：宋}
《宋史》为这一年留下了可回查的年序入口。

\noindent\textbf{宋。}\eventanchor{SYD-1030-NS-071}{八月丙戌，詔詳定鹽法}；\eventanchor{SYD-1030-NS-238}{丙申，弛三京、河中府、潁、許、汝、鄭、鄆、濟、衛、晉、絳、虢、亳、宿等二十八州軍鹽禁}。

\subsection{1031年：宋}
《宋史》为这一年留下了可回查的年序入口。

\noindent\textbf{宋。}\eventanchor{SYD-1031-NS-072}{丁酉，出知雜御史曹修古，御史郭勸、楊偕，推直官段少連}；\eventanchor{SYD-1031-NS-239}{冬十月丙戌，詔公卿大夫勵名節}。

\subsection{1032年：宋}
《宋史》为这一年留下了可回查的年序入口。

\noindent\textbf{宋。}\eventanchor{SYD-1032-NS-073}{丙申，皇太后出金銀器易左藏緡錢二十萬，以助修內}；\eventanchor{SYD-1032-NS-240}{丙午，詔仕廣南者毋過兩任，以防貪黷}。

\subsection{1033年：宋}
《宋史》为这一年留下了可回查的年序入口。

\noindent\textbf{宋。}\eventanchor{SYD-1033-NS-074}{丙辰，出道輔及諫官范仲淹，仍詔台諫自今毋相率請對}；\eventanchor{SYD-1033-NS-241}{丙辰，內侍江德明等並坐交通請謁黜}。

\subsection{1034年：宋}
《宋史》为这一年留下了可回查的年序入口。

\noindent\textbf{宋。}\eventanchor{SYD-1034-NS-075}{丙寅，詔開封府界諸縣作糜粥以濟饑民，諸災傷州軍亦如之}；\eventanchor{SYD-1034-NS-242}{帝為杖內侍，仍詔有司自今宮中傳命，毋得輒受}。

\subsection{1035年：宋}
《宋史》为这一年留下了可回查的年序入口。

\noindent\textbf{宋。}\eventanchor{SYD-1035-NS-076}{八月壬子朔，詔輕強盜法}；\eventanchor{SYD-1035-NS-243}{丙午，降天下系囚罪一等，杖以下釋之}。

\subsection{1036年：宋}
《宋史》为这一年留下了可回查的年序入口。

\noindent\textbf{宋。}\eventanchor{SYD-1036-NS-077}{八月己酉，班民間冠服、居室、車馬、器用犯制之禁}；\eventanchor{SYD-1036-NS-244}{冬十月丁未，命章得象等考課諸路提刑}。

\subsection{1037年：宋}
《宋史》为这一年留下了可回查的年序入口。

\noindent\textbf{宋。}\eventanchor{SYD-1037-NS-078}{八月甲戌，越州水，賜被溺民家錢有差}；\eventanchor{SYD-1037-NS-245}{冬十一月癸亥，罷登、萊賣金場}。

\subsection{1038年：宋}
《宋史》为这一年留下了可回查的年序入口。

\noindent\textbf{宋。}\eventanchor{SYD-1038-NS-079}{八月丁卯，復淮南、江、浙、荊湖制置發運使}；\eventanchor{SYD-1038-NS-246}{丙辰，以地震及雷發不時，詔轉運使、提舉刑獄按所部官吏}。

\subsection{1039年：宋}
《宋史》为这一年留下了可回查的年序入口。

\noindent\textbf{宋。}\eventanchor{SYD-1039-NS-080}{丙午，遣使體量安撫陝西、河東}；\eventanchor{SYD-1039-NS-247}{丙子，降三京囚罪一等，徒以下釋之}。

\subsection{1040年：宋}
《宋史》为这一年留下了可回查的年序入口。

\noindent\textbf{宋。}\eventanchor{SYD-1040-NS-081}{八月戊戌，禁以金箔飾佛像}；\eventanchor{SYD-1040-NS-248}{丙申，詔諸路轉運使、提刑訪知邊事者以聞}。

\subsection{1041年：宋}
《宋史》为这一年留下了可回查的年序入口。

\noindent\textbf{宋。}\eventanchor{SYD-1041-NS-082}{八月戊寅，詔鄜延部署以兵援麟、府}；\eventanchor{SYD-1041-NS-249}{丙辰，發廩粟，減價以濟京城民}。

\subsection{1042年：宋}
《宋史》为这一年留下了可回查的年序入口。

\noindent\textbf{宋。}\eventanchor{SYD-1042-NS-083}{丙辰，詔醫官毋得換右職}；\eventanchor{SYD-1042-NS-250}{丁未，詔軍校戰沒無子孫者，賜其家緡錢}。

\subsection{1043年：宋}
《宋史》为这一年留下了可回查的年序入口。

\noindent\textbf{宋。}\eventanchor{SYD-1043-NS-084}{八月乙未朔，命官詳定編敕}；\eventanchor{SYD-1043-NS-251}{丙辰，以春夏不雨，遣使祠禱于嶽瀆}。

\subsection{1044年：宋}
《宋史》为这一年留下了可回查的年序入口。

\noindent\textbf{宋。}\eventanchor{SYD-1044-NS-085}{八月辛卯，命賈昌朝領天下農田，范仲淹領刑法事}；\eventanchor{SYD-1044-NS-252}{保州雲翼軍殺官吏、據城叛}。

\subsection{1045年：宋}
《宋史》为这一年留下了可回查的年序入口。

\noindent\textbf{宋。}\eventanchor{SYD-1045-NS-086}{八月庚午，荊南府、岳州地震}；\eventanchor{SYD-1045-NS-253}{丙戌，杜衍罷，以賈昌朝同中書門下平章事兼樞密使、集賢殿大學士，王貽永為樞密使，宋庠參知政事，吳育、龐籍並為樞密副使}。

\subsection{1046年：宋}
《宋史》为这一年留下了可回查的年序入口。

\noindent\textbf{宋。}\eventanchor{SYD-1046-NS-087}{丙申，詔陝西市蕃部馬}；\eventanchor{SYD-1046-NS-254}{丙寅，以久旱，民多暍死，命京城增鑿井三百九十}。

\subsection{1047年：宋}
《宋史》为这一年留下了可回查的年序入口。

\noindent\textbf{宋。}\eventanchor{SYD-1047-NS-088}{壬子，禦正殿，復常膳}；\eventanchor{SYD-1047-NS-255}{八月乙丑，析河北為四路，各置都總管}。

\subsection{1048年：宋}
《宋史》为这一年留下了可回查的年序入口。

\noindent\textbf{宋。}\eventanchor{SYD-1048-NS-089}{壬戌，以霖雨，錄系囚}；\eventanchor{SYD-1048-NS-256}{又詔翰林學士、三司使、知開封府、御史中丞曰：「朕躬闕失，左右朋邪，中外險詐，州郡暴虐，法令有不便於民者，朕欲聞之，其悉以陳}。

\subsection{1049年：宋}
《宋史》为这一年留下了可回查的年序入口。

\noindent\textbf{宋。}\eventanchor{SYD-1049-NS-090}{冬十一月丙申，詔河北被災民八十以上及篤疾不能自存者，人賜米一石、酒一斗}；\eventanchor{SYD-1049-NS-257}{二月戊辰，以河北疫，遣使頒藥}。

\subsection{1050年：宋}
《宋史》为这一年留下了可回查的年序入口。

\noindent\textbf{宋。}\eventanchor{SYD-1050-NS-091}{丙寅，秀州地震，有聲如雷}；\eventanchor{SYD-1050-NS-258}{丁卯，詔中書門下省、兩制及太常官詳定大樂}。

\subsection{1051年：宋}
《宋史》为这一年留下了可回查的年序入口。

\noindent\textbf{宋。}\eventanchor{SYD-1051-NS-092}{丙辰，以孔氏子孫復知仙源縣事}；\eventanchor{SYD-1051-NS-259}{八月丙戌，遣使安撫京東、淮南、兩浙、荊湖、江南饑民}。

\subsection{1052年：宋}
《宋史》为这一年留下了可回查的年序入口。

\noindent\textbf{宋。}\eventanchor{SYD-1052-NS-093}{八月癸未，詔開封府比大風雨，民廬摧圮壓死者，官為祭斂之}；\eventanchor{SYD-1052-NS-260}{丙午，命余靖經制廣南盜賊事}。

\subsection{1053年：宋}
《宋史》为这一年留下了可回查的年序入口。

\noindent\textbf{宋。}\eventanchor{SYD-1053-NS-094}{八月丁酉朔，詔民訴災傷而監司不受者，聽州軍以狀聞}；\eventanchor{SYD-1053-NS-261}{丙戌，詔廣西都監蕭注等追捕智高}。

\subsection{1054年：宋}
《宋史》为这一年留下了可回查的年序入口。

\noindent\textbf{宋。}\eventanchor{SYD-1054-NS-095}{八月丁酉，詔：「前代帝王後嘗仕本朝，官八品以下，其祖父母、父母、妻子犯流以下罪，聽贖}；\eventanchor{SYD-1054-NS-262}{丙午，溫成皇后神主入廟}。

\subsection{1055年：宋}
《宋史》为这一年留下了可回查的年序入口。

\noindent\textbf{宋。}\eventanchor{SYD-1055-NS-096}{丙子，封孔子後為衍聖公}；\eventanchor{SYD-1055-NS-263}{丁酉，詔武臣有贓濫者毋得轉橫行，其立戰功者許之}。

\subsection{1056年：宋}
《宋史》为这一年留下了可回查的年序入口。

\noindent\textbf{宋。}\eventanchor{SYD-1056-NS-097}{丙戌，賜河北流民米，壓溺死者，賜其家錢有差}；\eventanchor{SYD-1056-NS-264}{庚子，賜致仕卿、監以上及曾任近侍之臣粟帛酒饌}。

\subsection{1057年：宋}
《宋史》为这一年留下了可回查的年序入口。

\noindent\textbf{宋。}\eventanchor{SYD-1057-NS-098}{辛亥，立內降關白二府法}；\eventanchor{SYD-1057-NS-265}{八月己酉，詔：每歲賜諸道節鎮、諸州錢有差}。

\subsection{1058年：宋}
《宋史》为这一年留下了可回查的年序入口。

\noindent\textbf{宋。}\eventanchor{SYD-1058-NS-099}{丁-\{丑\}-，詔裁定制科及進士高第人恩數}；\eventanchor{SYD-1058-NS-266}{丙辰，詔：「守令或貪恣耄昏，以弛為寬，以苛為察，以增賦斂為勞，以出入刑罰為能，而部使者莫之舉劾}。

\subsection{1059年：宋}
《宋史》为这一年留下了可回查的年序入口。

\noindent\textbf{宋。}\eventanchor{SYD-1059-NS-100}{庚子，詔內臣員多，權罷進養子入內}；\eventanchor{SYD-1059-NS-267}{丙子，復銀台司封駁制}。

\subsection{1060年：宋}
《宋史》为这一年留下了可回查的年序入口。

\noindent\textbf{宋。}\eventanchor{SYD-1060-NS-101}{丙申，詔待制、台諫官、正刺史以上各舉諸司使至三班使臣堪將領及行陣戰鬥者三人}；\eventanchor{SYD-1060-NS-268}{丙戌，命近臣同三司議均稅}。

\subsection{1061年：宋}
《宋史》为这一年留下了可回查的年序入口。

\noindent\textbf{宋。}\eventanchor{SYD-1061-NS-102}{戊寅，詔州縣長吏有清白不擾而實惠及民者，令本路監司保薦再任，政跡尤異，當加獎擢}；\eventanchor{SYD-1061-NS-269}{丙戌，詔淮南、江、浙水災，差官體量蠲稅}。

\subsection{1062年：宋}
《宋史》为这一年留下了可回查的年序入口。

\noindent\textbf{宋。}\eventanchor{SYD-1062-NS-103}{丙申，詔內藏庫、三司共出緡錢一百萬，助糴天下常平倉}；\eventanchor{SYD-1062-NS-270}{冬十月乙亥，皇子表辭所除官，賜詔不允}。

\subsection{1063年：宋}
《宋史》为这一年留下了可回查的年序入口。

\noindent\textbf{宋。}\eventanchor{SYD-1063-NS-104}{八年春正月辛亥，交阯貢馴象九}；\eventanchor{SYD-1063-NS-271}{癸亥，御內東門幄殿，優賜諸軍緡錢}。

\subsection{1064年：宋}
《宋史》为这一年留下了可回查的年序入口。

\noindent\textbf{宋。}\eventanchor{SYD-1064-NS-105}{丙辰，內侍都知任守忠坐不法，貶保信軍節度副使、蘄州安置}；\eventanchor{SYD-1064-NS-272}{丁巳，以上供米三萬石振宿、亳二州水災戶}。

\subsection{1065年：宋}
《宋史》为这一年留下了可回查的年序入口。

\noindent\textbf{宋。}\eventanchor{SYD-1065-NS-106}{帝書其後曰：「雨災專以戒朕不德，可更曰‘協德交修’}；\eventanchor{SYD-1065-NS-273}{己亥，以水災，罷開樂宴}。

\subsection{1066年：宋}
《宋史》为这一年留下了可回查的年序入口。

\noindent\textbf{宋。}\eventanchor{SYD-1066-NS-107}{丙寅，幸降聖院，謁神-\{御\}-殿}；\eventanchor{SYD-1066-NS-274}{賜文武官子爲父後者勳一轉}。

\subsection{1067年：宋}
《宋史》为这一年留下了可回查的年序入口。

\noindent\textbf{宋。}\eventanchor{SYD-1067-NS-108}{一日，語神宗曰：「國家舊制，士大夫之子有尚帝女，皆升行以避舅姑之尊，義甚無謂}；\eventanchor{SYD-1067-NS-275}{初辭皇子，請潭王宮教授周孟陽作奏，孟陽有所勸戒，即謝而拜之}。

\subsection{1068年：宋}
《宋史》为这一年留下了可回查的年序入口。

\noindent\textbf{宋。}\eventanchor{SYD-1068-NS-109}{八月壬寅，詔京東、西路存恤河北流民}；\eventanchor{SYD-1068-NS-276}{丙申，趙概罷知徐州，三司使唐介參知政事}。

\subsection{1069年：宋}
《宋史》为这一年留下了可回查的年序入口。

\noindent\textbf{宋。}\eventanchor{SYD-1069-NS-110}{丙辰，詔-\{御\}-史請對，並許直由閣門上殿}；\eventanchor{SYD-1069-NS-277}{丙午，同修起居注範純仁以言事多忤安石，罷同知諫院}。

\subsection{1070年：宋}
《宋史》为这一年留下了可回查的年序入口。

\noindent\textbf{宋。}\eventanchor{SYD-1070-NS-111}{八月戊午，罷看詳銀台司文字所}；\eventanchor{SYD-1070-NS-278}{丙辰，立試刑法及詳刑官}。

\subsection{1071年：宋}
《宋史》为这一年留下了可回查的年序入口。

\noindent\textbf{宋。}\eventanchor{SYD-1071-NS-112}{八月癸-\{丑\}-朔，高麗來貢}；\eventanchor{SYD-1071-NS-279}{丙辰，置樞密院檢詳官}。

\subsection{1072年：宋}
《宋史》为这一年留下了可回查的年序入口。

\noindent\textbf{宋。}\eventanchor{SYD-1072-NS-113}{八月甲申，太子少師致仕歐陽修薨}；\eventanchor{SYD-1072-NS-280}{丙午，以內藏庫錢置市易務}。

\subsection{1073年：宋}
《宋史》为这一年留下了可回查的年序入口。

\noindent\textbf{宋。}\eventanchor{SYD-1073-NS-114}{丙辰，以四月朔日當食，自丁巳避殿、減膳，降天下囚罪一等，流以下釋之}；\eventanchor{SYD-1073-NS-281}{丙申，永昌陵上宮東門火}。

\subsection{1074年：宋}
《宋史》为这一年留下了可回查的年序入口。

\noindent\textbf{宋。}\eventanchor{SYD-1074-NS-115}{八月丁-\{丑\}-，賜環慶安撫司度僧牒，以募粟振漢蕃饑民}；\eventanchor{SYD-1074-NS-282}{丙午，遣使分行諸路，募武士赴熙河}。

\subsection{1075年：宋}
《宋史》为这一年留下了可回查的年序入口。

\noindent\textbf{宋。}\eventanchor{SYD-1075-NS-116}{八年春正月庚子，蔡挺罷判南京留司-\{御\}-史台，馮京罷知亳州}；\eventanchor{SYD-1075-NS-283}{丙辰，-\{御\}-殿復膳}。

\subsection{1076年：宋}
《宋史》为这一年留下了可回查的年序入口。

\noindent\textbf{宋。}\eventanchor{SYD-1076-NS-117}{八月甲申朔，齊州監務左班殿直孫紀死賊，錄其一子爲三班借職}；\eventanchor{SYD-1076-NS-284}{白丁王禹錫等二人，賜錢其家}。

\subsection{1077年：宋}
《宋史》为这一年留下了可回查的年序入口。

\noindent\textbf{宋。}\eventanchor{SYD-1077-NS-118}{辛酉，詔鎮戎、德順軍各置都監一員}；\eventanchor{SYD-1077-NS-285}{乙卯，詔：「諸傳宣、內批、面諭，事無法守，並從中書、樞密覆奏}。

\subsection{1078年：宋}
《宋史》为这一年留下了可回查的年序入口。

\noindent\textbf{宋。}\eventanchor{SYD-1078-NS-119}{丙辰，詔青州民王贇以復父仇免死，刺配鄰州}；\eventanchor{SYD-1078-NS-286}{丙辰，詔增置兩浙路提舉官}。

\subsection{1079年：宋}
《宋史》为这一年留下了可回查的年序入口。

\noindent\textbf{宋。}\eventanchor{SYD-1079-NS-120}{丙午，復置-\{御\}-史六察}；\eventanchor{SYD-1079-NS-287}{丙戌，詔雄州兩輸戶南徙者諭令復業}。

\subsection{1080年：宋}
《宋史》为这一年留下了可回查的年序入口。

\noindent\textbf{宋。}\eventanchor{SYD-1080-NS-121}{八月乙巳，罷省、寺、監官領空名者}；\eventanchor{SYD-1080-NS-288}{丙辰，始-\{御\}-崇政殿視朝}。

\subsection{1081年：宋}
《宋史》为这一年留下了可回查的年序入口。

\noindent\textbf{宋。}\eventanchor{SYD-1081-NS-122}{八月乙卯朔，罷中書堂選，悉歸有司}；\eventanchor{SYD-1081-NS-289}{丙辰，詔蠲河北東路災傷州軍今年夏料役錢}。

\subsection{1082年：宋}
《宋史》为这一年留下了可回查的年序入口。

\noindent\textbf{宋。}\eventanchor{SYD-1082-NS-123}{八月庚戌朔，封-\{御\}-侍武氏爲才人}；\eventanchor{SYD-1082-NS-290}{丙辰，以乞弟平，班師}。

\subsection{1083年：宋}
《宋史》为这一年留下了可回查的年序入口。

\noindent\textbf{宋。}\eventanchor{SYD-1083-NS-124}{八月丙子，賜升祔陪祠官宴於尚書省}。

\subsection{1084年：宋}
《宋史》为这一年留下了可回查的年序入口。

\noindent\textbf{宋。}\eventanchor{SYD-1084-NS-125}{八月庚午，詔王光祖遣人招諭乞弟，許出降免罪補官}。

\subsection{1085年：宋}
《宋史》为这一年留下了可回查的年序入口。

\noindent\textbf{宋。}\eventanchor{SYD-1085-NS-126}{丁亥，命禮部鎖試別所}。

\subsection{1086年：宋}
《宋史》为这一年留下了可回查的年序入口。

\noindent\textbf{宋。}\eventanchor{SYD-1086-NS-127}{八月辛卯，詔常平依舊法，罷青苗錢}。

\subsection{1087年：宋}
《宋史》为这一年留下了可回查的年序入口。

\noindent\textbf{宋。}\eventanchor{SYD-1087-NS-128}{八月辛巳，程頤罷經筵，權同管勾西京國子監}。

\subsection{1088年：宋}
《宋史》为这一年留下了可回查的年序入口。

\noindent\textbf{宋。}\eventanchor{SYD-1088-NS-129}{八月戊寅，阿-\{里\}-骨入貢}。

\subsection{1089年：宋}
《宋史》为这一年留下了可回查的年序入口。

\noindent\textbf{宋。}\eventanchor{SYD-1089-NS-130}{丁亥，復貶確爲英州別駕、安置新州}。

\subsection{1090年：宋}
《宋史》为这一年留下了可回查的年序入口。

\noindent\textbf{宋。}\eventanchor{SYD-1090-NS-131}{丙辰，禁軍大閱，賜以銀楪、匹帛，罷轉資}。

\subsection{1091年：宋}
《宋史》为这一年留下了可回查的年序入口。

\noindent\textbf{宋。}\eventanchor{SYD-1091-NS-132}{八月己-\{丑\}-，三省進納後六禮儀制}。

\subsection{1092年：宋}
《宋史》为这一年留下了可回查的年序入口。

\noindent\textbf{宋。}\eventanchor{SYD-1092-NS-133}{己酉，永興軍、蘭州、鎮戎軍地震}。

\subsection{1093年：宋}
《宋史》为这一年留下了可回查的年序入口。

\noindent\textbf{宋。}\eventanchor{SYD-1093-NS-134}{壬子，詔刑部不得分禁系人數，瘐死數多者申尚書省}。

\subsection{1094年：宋}
《宋史》为这一年留下了可回查的年序入口。

\noindent\textbf{宋。}\eventanchor{SYD-1094-NS-135}{編類元祐群臣章疏及更改事條}。

\subsection{1095年：宋}
《宋史》为这一年留下了可回查的年序入口。

\noindent\textbf{宋。}\eventanchor{SYD-1095-NS-136}{八月壬申，命彰信軍節度使宗景爲開府儀同三司，封濟陰郡王}。

\subsection{1096年：宋}
《宋史》为这一年留下了可回查的年序入口。

\noindent\textbf{宋。}\eventanchor{SYD-1096-NS-137}{丙戌，詔三歲一取旨，遣郎官、-\{御\}-史按察監司職事}。

\subsection{1097年：宋}
《宋史》为这一年留下了可回查的年序入口。

\noindent\textbf{宋。}\eventanchor{SYD-1097-NS-138}{八月乙酉，封湖州觀察使世開爲安定郡王}。

\subsection{1098年：宋}
《宋史》为这一年留下了可回查的年序入口。

\noindent\textbf{宋。}\eventanchor{SYD-1098-NS-139}{八月丙子朔，熙河蘭岷路復爲熙河蘭會路}。

\subsection{1099年：宋}
《宋史》为这一年留下了可回查的年序入口。

\noindent\textbf{宋。}\eventanchor{SYD-1099-NS-140}{八月癸酉，章惇等進《新修敕令式》}。

\subsection{1100年：宋}
《宋史》为这一年留下了可回查的年序入口。

\noindent\textbf{宋。}\eventanchor{SYD-1100-NS-141}{皇太后諭遺制，立弟端王即位於柩前，皇太后權同處分軍國事}。

\subsection{1101年：宋}
《宋史》为这一年留下了可回查的年序入口。

\noindent\textbf{宋。}\eventanchor{SYD-1101-NS-142}{丁丑，易大行皇太后園陵為山陵，命曾布為山陵使}。

\subsection{1102年：宋}
《宋史》为这一年留下了可回查的年序入口。

\noindent\textbf{宋。}\eventanchor{SYD-1102-NS-143}{八月乙卯，子烜改名桓，煥改名楷}。

\subsection{1103年：宋}
《宋史》为这一年留下了可回查的年序入口。

\noindent\textbf{宋。}\eventanchor{SYD-1103-NS-144}{丙寅，詔六曹長貳歲考郎官治狀，分三等以聞}。

\subsection{1104年：宋}
《宋史》为这一年留下了可回查的年序入口。

\noindent\textbf{宋。}\eventanchor{SYD-1104-NS-145}{八月庚子，詔諸路知州、通判增入「主管學事」四字}。

\subsection{1105年：宋}
《宋史》为这一年留下了可回查的年序入口。

\noindent\textbf{宋。}\eventanchor{SYD-1105-NS-146}{丙辰，詔自今非宰臣毋得除特進}。

\subsection{1106年：宋}
《宋史》为这一年留下了可回查的年序入口。

\noindent\textbf{宋。}\eventanchor{SYD-1106-NS-147}{丙寅，蔡京罷為開府儀同三司、中太一宮使}。

\subsection{1107年：宋}
《宋史》为这一年留下了可回查的年序入口。

\noindent\textbf{宋。}\eventanchor{SYD-1107-NS-148}{朝散郎吳儲、承議郎吳侔坐與妖人張懷素謀反，伏誅}。

\subsection{1108年：宋}
《宋史》为这一年留下了可回查的年序入口。

\noindent\textbf{宋。}\eventanchor{SYD-1108-NS-149}{八月辛巳，邢州河水溢，壞民廬舍，復被水者家}。

\subsection{1109年：宋}
《宋史》为这一年留下了可回查的年序入口。

\noindent\textbf{宋。}\eventanchor{SYD-1109-NS-150}{八月乙酉，封子朴為雍國公}。

\subsection{1110年：宋}
《宋史》为这一年留下了可回查的年序入口。

\noindent\textbf{宋。}\eventanchor{SYD-1110-NS-151}{丙辰，詔以彗見，避殿減膳，令侍從官直言指陳闕失}。

\subsection{1111年：宋}
《宋史》为这一年留下了可回查的年序入口。

\noindent\textbf{宋。}\eventanchor{SYD-1111-NS-152}{八月乙未，復蔡京為太子太師}。

\subsection{1112年：宋}
《宋史》为这一年留下了可回查的年序入口。

\noindent\textbf{宋。}\eventanchor{SYD-1112-NS-153}{成都路夷人董舜咨、董彥博內附，置祺、亨二州}。

\subsection{1113年：宋}
《宋史》为这一年留下了可回查的年序入口。

\noindent\textbf{宋。}\eventanchor{SYD-1113-NS-154}{丙申，升蘇州為平江府}。

\subsection{1114年：宋}
《宋史》为这一年留下了可回查的年序入口。

\noindent\textbf{宋。}\eventanchor{SYD-1114-NS-155}{八月乙巳，改端明殿學士為延康殿學士，樞密直學士為述古殿直學士}。

\subsection{1115年：宋}
《宋史》为这一年留下了可回查的年序入口。

\noindent\textbf{宋。}\eventanchor{SYD-1115-NS-156}{八月己酉，以秘書省地為明堂}。

\subsection{1116年：宋}
《宋史》为这一年留下了可回查的年序入口。

\noindent\textbf{宋。}\eventanchor{SYD-1116-NS-157}{丁丑，詔天甯諸節及壬戌日，杖已下罪聽贖}。

\subsection{1117年：宋}
《宋史》为这一年留下了可回查的年序入口。

\noindent\textbf{宋。}\eventanchor{SYD-1117-NS-158}{八月癸亥，詔明堂並祠五帝}。

\subsection{1118年：宋}
《宋史》为这一年留下了可回查的年序入口。

\noindent\textbf{宋。}\eventanchor{SYD-1118-NS-159}{丙戌，詔太學、辟雍各置《內經》、《道德經》、《莊子》、《列子》博士二員}。

\subsection{1119年：宋}
《宋史》为这一年留下了可回查的年序入口。

\noindent\textbf{宋。}\eventanchor{SYD-1119-NS-160}{改女冠為女道，尼為女德}。

\subsection{1120年：宋}
《宋史》为这一年留下了可回查的年序入口。

\noindent\textbf{宋。}\eventanchor{SYD-1120-NS-161}{八月庚辰，詔減定醫官額}。

\subsection{1121年：宋}
《宋史》为这一年留下了可回查的年序入口。

\noindent\textbf{宋。}\eventanchor{SYD-1121-NS-162}{陳過庭、張汝霖以乞罷御前使喚及歲進花果，為王黼所劾，並竄貶}。

\subsection{1122年：宋}
《宋史》为这一年留下了可回查的年序入口。

\noindent\textbf{宋。}\eventanchor{SYD-1122-NS-163}{丙申，貶劉延慶為率府率、安置筠州}。

\subsection{1123年：宋}
《宋史》为这一年留下了可回查的年序入口。

\noindent\textbf{宋。}\eventanchor{SYD-1123-NS-164}{丁酉，進封雍國公朴為華原郡王，徐國公棣為高平郡王，並為開府儀同三司}。

\subsection{1124年：宋}
《宋史》为这一年留下了可回查的年序入口。

\noindent\textbf{宋。}\eventanchor{SYD-1124-NS-165}{八月乙卯，譚稹落太尉、罷宣撫使，童貫落致仕，領樞密院代之}。

\subsection{1125年：宋}
《宋史》为这一年留下了可回查的年序入口。

\noindent\textbf{宋。}\eventanchor{SYD-1125-NS-166}{罷道官，罷大晟府、行幸局}。

\subsection{1126年：宋}
《宋史》为这一年留下了可回查的年序入口。

\noindent\textbf{宋。}\eventanchor{SYD-1126-NS-167}{內侍朱拱之宣綱後期，眾臠而磔之，並殺內侍數十人}。
